\scanpage{52}
% archive sheet number: 51
\mspage{2 - 7 -}\hfill n\textsuperscript{o} 392

Pour tout $x\handwritten{\in}K$, soit $\handwritten{\Delta}(x)$ l'élément $(x_{\unsure{\underset{=}{p}|\underline{p}|p}})\handwritten{\in\prod}_{\underset{=}{p}\handwritten{\in}P}\overset{\handwritten{\scriptscriptstyle\wedge}}{K}_{\underset{=}{p}}$ tel que $x_{\underset{=}{p}}=x$ pour tout $\underset{=}{p}\handwritten{\in}P$ ; comme $x_{\underset{=}{p}}\handwritten{\in}\overset{\handwritten{\scriptscriptstyle\wedge}}{A}_{\underset{=}{p}}$ sauf pour un nombre fini de valeurs de $\underset{=}{p}$, on a $\handwritten{\Delta}(x)\handwritten{\in}\underline{A}$ ; on définit donc ainsi un homomorphisme $\handwritten{\Delta}:K\longrightarrow \underline{A}$, qui est \uline{injectif} si $P\handwritten{\ne\varnothing}$, (c'est-à-dire si $A$ n'est pas un corps) ; les éléments de $\handwritten{\Delta}(K)$ sont dits \uline{adèles restreints principaux}, et il est clair que $\handwritten{\Delta}(A)\handwritten{\subset}\underline{A}_{\unsure{0|o}}$.

PROPOSITION 3.- \uline{L'anneau} $\underline{A}_0$, (resp. $\underline{A}$), \uline{s'identifie au complété de} $A$, (resp. $K$), \uline{pour la topologie d'anneau dont un système fondamental de voisinages de} $0$ \uline{est formé de tous les idéaux entiers} $\handwritten{\ne}(0)$ \uline{de} $A$.

Il est immédiat que la topologie considérée sur $A$, (ou $K$), est séparée. Compte tenu du n\textsuperscript{o} 3, l'assertion relative à $\underline{A}_0$ résulte du chap. III, § 2, n\textsuperscript{o} 13, prop. 17. Cela montre donc que $\handwritten{\Delta}(A)$ est dense dans $\underline{A}_0$ ; pour voir de même que $\handwritten{\Delta}(K)$ est dense dans $\underline{A}$, on remarque que pour tout $\underline{x}=(x_{\underset{=}{p}})\handwritten{\in}\underline{A}$, il n'y a qu'un nombre fini de $\underset{=}{p}\handwritten{\in}P$ tels que $v_{\underset{=}{p}}(x_{\underset{=}{p}})\handwritten{<}0$ ; en vertu du § 1, n\textsuperscript{o} 5, prop. 9, il y a donc un $s\handwritten{\in}K$ tel que $sx_{\underset{=}{p}}\handwritten{\in}\overset{\handwritten{\scriptscriptstyle\wedge}}{A}_{\underset{=}{p}}$ pour tout $\underset{=}{p}\handwritten{\in}P$, autrement dit $\handwritten{\Delta}(s)\underline{x}\handwritten{\in}\underline{A}_0$, et comme la multiplication par $\handwritten{\Delta}(s)$ est un homéomorphisme de $\underline{A}$ sur lui-même, il suffit d'appliquer le fait que $\handwritten{\Delta}(A)$ est dense dans $\underline{A}_0$ pour en déduire que $\handwritten{\Delta}(K)$ est dense dans $\underline{A}$.

On pourrait naturellement aussi prouver que $\handwritten{\Delta}(K)$ est dense dans $\underline{A}$ en utilisant la prop. 2.

Considérons maintenant le groupe multiplicatif $\mathrm{SL}(n,\underline{A})$, formé des matrices $\underline{U}\handwritten{\in}M_n(\underline{A})$ telles que $\det(\underline{U})=1$ ; si on munit $M_n(\underline{A})=\underline{A}^{n^2}$ de la topologie produit, elle induit sur $\mathrm{SL}(n,\underline{A})$ une topologie \uline{compatible avec la structure de groupe} de $\mathrm{SL}(n,\underline{A})$. En effet, il suffit de voir que l'application $\underline{U}\longrightarrow \underline{U}^{-1}$ est continue dans $\mathrm{SL}(n,\underline{A})$ ; mais comme $\underline{U}$ est unimodulaire, on sait, (\uline{Alg}., chap. III, 3e éd., § \hspace{1em}, n\textsuperscript{o} \hspace{1em}), que les éléments de $\underline{U}^{-1}$ sont des \uline{mineurs} de $\underline{U}$, donc des polynômes en les éléments de $\underline{U}$, ce qui prouve notre assertion. Si on identifie $K$ à un sous-anneau
